\clearpage
\SourceLeaf{207}
\renewcommand{\BookRunningTitle}{Manuscript Papers}
\thispagestyle{empty}
\pdfbookmark[0]{Manuscript Papers}{handschriftlicher-nachlass}
\vspace*{.25\textheight}
\begin{center}
{\Large MANUSCRIPT}\par\vspace{2.8cm}
{\Huge PAPERS.}
\end{center}
\vfill

\clearpage
\SourceLeaf{209}
\renewcommand{\BookRunningTitle}{Manuscript Papers. Analysis of Residues}
\pdfbookmark[0]{Solution of a Pure Congruence. Analysis of Residues}{analysis-residuorum-pars-prior}
\begin{center}
{\large SOLUTION OF THE CONGRUENCE $X^m-1\equiv0$.}\par\medskip
ANALYSIS OF RESIDUES. CHAPTER SIX. FIRST PART.
\end{center}
\BookRule{.18\linewidth}
\renewcommand{\Article}[1]{\par\medskip\Needspace{4\baselineskip}\pdfbookmark[1]{Art. #1}{analysis-residuorum-art-#1}\centerline{\textbf{#1.}}\nobreak\smallskip}
\Article{237}
In Chapter \textsc{iii} we showed that the congruence $x^n\equiv1$, when a prime number $p$ is taken as modulus, has $\mu$ roots, where $\mu$ is the greatest common divisor of the numbers $n$ and $p-1$, and that these roots agree exactly with the roots of the congruence $x^\mu\equiv1$. It therefore suffices to consider the case in which $n$ is a divisor of $p-1$. That the solution not only of the congruence $x^n-1\equiv0$ but of any other congruence for arbitrary moduli can be derived from its solution for prime moduli has already been shown at various places and will be developed more fully below (Chapter \textsc{viii}).

\Article{238}
But there is no need to stop even here. In the same Chapter \textsc{iii} we explained that the solution of the congruence $x^n\equiv1$ depends on the solution of similar congruences $x^a\equiv1$, $x^b\equiv1$, and so on, where $a,b$, and so on, are primes or powers of primes, and $n$ is the product of these numbers. Namely, if $A,B$, and so on, are respectively any roots of the congruences $x^a\equiv1$, $x^b\equiv1$, and so on, their product $AB\ldots$ will be one of the roots of the congruence $x^n\equiv1$. Our investigations will therefore be confined to solving the congruence $x^n\equiv1\;(\mathrm{mod.}\,p)$, where $p$ is a prime, and $n$ a prime or a power of a prime and also a divisor of $p-1$.

\SourceLeaf{210}
\Article{239}
Furthermore, Chapter \textsc{iii} shows that among the roots of the congruence $x^n\equiv1$ there are always some whose powers can represent all the others. Thus, if $r$ denotes a root of this kind (we called it \emph{primitive} above when $n=p-1$, and shall retain this expression here, though in a broader sense), all the roots of the proposed congruence will be
\[
1,\ r,\ rr,\ r^3\quad\dotsfive\quad r^{n-1}
\]
We must therefore direct every effort towards finding roots of this kind, since once they have been found the others will become evident of themselves. For brevity we denote any power of $r$ by its exponent enclosed in parentheses, so that $(0)$ denotes unity, $(1)$ any primitive root of the congruence $x^n\equiv1$, $(2)$ the square of $(1)$, and so on. Thus the sequence $(0),\ (1),\ (2),\ (3),\dotsfour(n-1)$ includes all the roots. It is also known that $(k)$ will always be such a primitive root whenever $k$ is coprime to $n$; that is, in our case, where $n$ is a power of the prime $t$, $=t^\nu$, whenever $t$ does not divide $k$. The symbols $(1)$, $(2)$, and so on, are evidently undetermined in themselves; but as soon as a definite value is assigned to $(1)$, all the others become determined.

\Article{240}
Since our aim is to investigate primitive roots before the others, we must first separate them from the rest. This is done by removing from the sequence $(0),(1),(2)\ldots(n-1)$ all terms $(k)$ for which $k$ is divisible by $t$. If $n$ is prime, or $\nu=1$, only $(0)$ is to be removed. Before proceeding to investigate the remaining roots, however, we earnestly advise the reader to work out several examples, so as to see concretely everything that might otherwise seem too general. We give one example; but this will not make it unnecessary to work out others independently.

Let $p=29$; $n=7$. The seven roots of the congruence $x^7\equiv1\;(\mathrm{mod.}\,29)$ will be $1,7,16,20,23,24,25$. Since $n$ is prime, all these roots except $1$ will be primitive. Thus, if $7=(1)$, the symbols will have the following meanings:
\[
\begin{array}{ccccccc}
(0)&(1)&(2)&(3)&(4)&(5)&(6)\\
1&7&20&24&23&16&25
\end{array}
\]

\SourceLeaf{211}
Everyone will also remember that the symbols $(n)$ and $(0)$, $(n+1)$ and $(1)$, and so on, and in general $(a)$ and $(b)$, are equivalent whenever $a\equiv b\;(\mathrm{mod.}\,n)$.

\Article{241}
For our purpose, however, we must proceed in yet another way. We retain only those terms $(k)$ for which $k$ is not divisible by $t$, whose number is $\frac{t-1}{t}\pdot n=\lambda$. All these numbers, or numbers congruent to them modulo $n$, can be represented by successive powers of some number. Let this be $=\rho$; then all the primitive roots of the congruence $x^n\equiv1$ will be denoted as follows:
\[
(1)\quad(\rho)\quad(\rho^2)\quad(\rho^3)\ \dotsfour\ (\rho^{\lambda-1})
\]
By this device we completely exclude all nonprimitive roots; the reasons and advantages will become clearer below. In our example we may therefore put $\rho=3$, and the primitive roots of the congruence $x^7\equiv1$ are arranged as follows:
\[
\begin{array}{rcccccc}
&(1)&(3)&(3^2)&(3^3)&(3^4)&(3^5)\\
\text{or}&(1)&(3)&(2)&(6)&(4)&(5)\\
\text{which are}&7&24&20&25&23&16
\end{array}
\]

\Needspace{10\baselineskip}\Article{242}
To let the reader know the aim of the following investigations, it will be useful to state the theorem that we propose to prove and clarify.

\emph{If the number $\lambda$, which is $=t^{\nu-1}\pdot t-1$, has prime factors $a,b,c,d$, and so on, and $\lambda=a^\alpha b^\beta c^\gamma\ldots$, the solution of the congruence $x^n-1\equiv0$ depends on solving $\alpha+\beta+\dotstwo$ congruences of lower degree, of which $\alpha$ have degree $a$, $\beta$ degree $b$, $\gamma$ degree $c$, and so on.}

Thus, in our example, solving the congruence $x^7\equiv1$ depends on a congruence of degree two and another of degree three; and it is apparent in general that the degrees of these congruences never depend on the modulus $p$. To reach the proof of this theorem, however, we must first give some propositions concerning the connection between congruences and their roots, although these investigations properly belong to Chapter eight, where they are to be pursued further.

\SourceLeaf{212}
\Article{243}
\textsc{Theorem.}\quad\emph{Suppose the congruence}
\[
x^m+Ax^{m-1}+Bx^{m-2}+\dotsfour+N\equiv0\;(\text{prime modulus})
\]
\emph{is such that, when the product of the $m$ factors $x-r$, $x-r'$, $x-r''$, $x-r'''\ldots$ is expanded to give $x^m+ax^{m-1}+bx^{m-2}\ldots+n$, we have $A\equiv a$, $B\equiv b$, $C\equiv c$, and so on, modulo $p$. Then the quantities $r,r',r''\ldots$ will be roots of the proposed congruence, and it will have no others.}

\emph{Proof.}\quad I. We always have
\[
\begin{aligned}
&x^m+Ax^{m-1}+Bx^{m-2}+\ldots\\
&\qquad\equiv x^m+ax^{m-1}+bx^{m-2}+\ldots\;(\mathrm{mod.}\,p)
\end{aligned}
\]
But the latter side of the congruence becomes $=0$ on setting $x=r$, $x=r'$, $x=r''$, and so on; therefore, for these values of $x$, the former side becomes $\equiv0\;(\mathrm{mod.}\,p)$. This proves the first part.

II. If another value $\rho$, congruent to none of $r,r'$, and so on, also satisfied the proposed congruence, we would have
\[
\begin{aligned}
0&\equiv\rho^m+A\rho^{m-1}+B\rho^{m-2}+\ldots\\
&\equiv\rho^m+a\rho^{m-1}+b\rho^{m-2}+\ldots\\
&\equiv(\rho-r)(\rho-r')(\rho-r'')(\rho-r''')\ldots
\end{aligned}
\]
But since none of the factors $\rho-r$, $\rho-r'$, $\rho-r''$, and so on, is $\equiv0$, their product being $\equiv0$ is impossible, because $p$ is prime. There are therefore no roots besides $r,r'$, and so on. This proves the second part.

\Needspace{12\baselineskip}\Article{244}
\textsc{Problem.}\quad\emph{Let $r,r',r''\ldots$ be unknown quantities, whose number is $=m$, whose sum is $=\alpha$, sum of squares $=\beta$, sum of cubes $=\gamma\ldots$, and sum of powers with exponent $m$ is $=\mu$. Suppose that these numbers, also $=m$ in number, are not themselves given, but instead other numbers $\alpha',\beta',\gamma'$, and so on, congruent to them individually modulo $p$, where this is a prime and $>m$. Find a congruence of degree $m$ whose roots are $r,r',r''$, and so on.}

\emph{Solution.}\quad Regard $r,r',r''$, and so on, as the roots of an equation
\[
x^m+Ax^{m-1}+Bx^{m-2}+Cx^{m-3}+\ldots=0
\]
and determine its coefficients $A,B,C$, and so on, by the known method, using congruence in place of equality; namely, put
\SourceLeaf{213}
\[
\begin{aligned}
-A&\equiv\alpha'\\
-2B&\equiv\beta'+A\alpha'\\
-3C&\equiv\gamma'+A\beta'+B\alpha'\\
-4D&\equiv\delta'+A\gamma'+B\beta'+C\alpha'\\
&\text{etc.}\\
-mN&\equiv\mu'+A\lambda'+\text{etc.}
\end{aligned}
\]
\begingroup\rightskip=0pt plus2em\relax These coefficients cannot be undetermined, since all the numbers \mbox{$1,2,3\ldots m<p$}. I say that the congruence
\[
x^m+Ax^{m-1}+Bx^{m-2}+\ldots+N\equiv0
\]
is the one required.\par\endgroup

\emph{Proof.}\quad Let the equation whose roots are $r,r',r'',r'''$, and so on, be
\[
x^m+ax^{m-1}+bx^{m-2}+\ldots=0
\]
Then
\[
\begin{aligned}
-a&=\alpha\\
-2b&=\beta+a\alpha\\
-3c&=\gamma+a\beta+b\alpha\\
-4d&=\delta+a\gamma+b\beta+c\alpha\\
&\text{etc.}
\end{aligned}
\]
\Needspace{12\baselineskip}It will be evident to everyone from this that
\[
a\equiv A,\quad b\equiv B,\quad c\equiv C\ \text{etc.}\;(\mathrm{mod.}\,p)
\]
so, by the preceding section, the numbers $r,r',r''$, and so on, which are roots of the \emph{equation}
\[
x^m+ax^{m-1}+bx^{m-2}+\ldots=0
\]
will also be roots of the \emph{congruence}
\[
x^m+Ax^{m-1}+Bx^{m-2}+\ldots\equiv0\qquad\text{Q. E. D.}
\]
We leave the construction of examples to the readers.

\Article{245}
We return to our purpose. Retaining the symbols used in sections 242 and earlier, we undertake to show that, if $\lambda$ is the product of any factors
\SourceLeaf{214}
$efg$, and so on, the primitive roots of the congruence $x^n\equiv1$, whose number is $\lambda$, can be divided into $e$ classes in such a way that the sums of the roots assigned to each class are given by a congruence of degree $e$. Once these are taken as known, each class can be subdivided into $f$ groups so that the sum of each group is given by a congruence of degree $f$. These groups can be subdivided again, and so on, until the individual roots are reached.

\Article{246}
\emph{Definition.}\quad We shall call the collection of \emph{all} terms contained in a form $(\rho^{ke+\alpha})$ of this kind (section 241) a \emph{complete period}, or simply a \emph{period}. Here $e$ denotes any divisor of $\lambda$; $\alpha$ any given number; and $k$ all integers from $0$ to $\frac{\lambda}{e}-1$. For brevity we denote such a period by $(e\ast\alpha)$. Thus, in our example, the terms
\[
\begin{array}{rclcl}
&(1),(2),(4)&\text{form the periods}&(2\ast0)&\text{,}\\
&(3),(6),(5)&&(2\ast1)&\\
\text{and}&(1),(6)&\text{form these}&(3\ast0)&\\
&(3),(4)&&(3\ast1)&\\
&(2),(5)&&(3\ast2)&
\end{array}
\]

\emph{Now, if all terms are distributed into periods in any way, and each period is again divided into smaller periods, and so on, we say that what we promised in the preceding section is achieved.}

Before entering into this exposition itself, however, we shall show that, although the formation of such a period depends on the two somewhat arbitrary quantities $r,\rho$, it contains nothing indeterminate: however these quantities are chosen, the same terms always come together in the same period, provided that the number of terms the period is to contain has been prescribed.

The criterion that two terms $A,B$ belong to the same period is that each is contained in a form $(\rho^{ke+\alpha})$ of this kind, or that $A\equiv r^{\rho^{ke+\alpha}}$, $B\equiv r^{\rho^{k'e+\alpha}}\;(\mathrm{mod.}\,p)$. Here $r$ is a primitive root of the congruence $x^n\equiv1\;(\mathrm{mod.}\,p)$, and $\rho$ a primitive root of the congruence $x^\lambda\equiv1\;(\mathrm{mod.}\,n)$; see above.

We must show that, if, in place of the numbers $r,\rho$, other numbers, say $s,\sigma$, are chosen, then $A$ and $B$ are contained in similar forms $s^{\sigma^{le+\beta}}$, $s^{\sigma^{l'e+\beta}}$.

\SourceLeaf{215}
Let $s^m\equiv r\;(\mathrm{mod.}\,p)$; $\sigma^\mu\equiv\rho\;(\mathrm{mod.}\,n)$ and $m\equiv\sigma^\zeta\;(\mathrm{mod.}\,n)$, which is possible because $r,\rho$ are primitive roots. Moreover, $m$ will be coprime to $n$, and $\mu$ to $\lambda$ (Chapter \textsc{iii}). By the appropriate substitutions we obtain
\[
A\equiv s^{\sigma^{\mu ke+\mu\alpha+\zeta}},\qquad B=s^{\sigma^{\mu k'e+\mu\alpha+\zeta}}\qquad\text{Q. E. D.}
\]
